\section{The nine modules that close the manuscript contract}
\subsection{Row 143: actual pointwise projectors}
The earlier coefficient projector and formal Fourier average are joined to the literal finite average on functions.
\excerpt{PointwiseFilters}{root_average}
The projector law is proved for actual pointwise values:
\excerpt{PointwiseFilters}{pointwise_projector}
On an invariant domain, the range characterization also tracks differentiability and equality restricted to that domain.
\excerpt{PointwiseFilters}{pointwise_range}
The domain condition says multiplication by the primitive root maps the set into itself. This is needed to apply covariance repeatedly along an orbit without evaluating a domain-restricted hypothesis outside its scope. The proof of the reverse implication chooses the function itself as a preimage. The forward implication uses covariance of the finite root average. Row 143 also includes the full pointwise product law and coefficient-series summation, closing the formal-to-holomorphic clauses of results 3.1--3.3.

\subsection{Row 144: arbitrary-argument addition}
\code{PointwiseAddition} supplies the actual entire transfer addition theorem and the signed addition theorem for arbitrary complex arguments. The signed carry is derived from the relation $\rho^N=-1$. This completes results 4.4 and 5.5 together with their earlier formal rows. No unspecified analytic evaluation map is left between the bivariate formal identity and the pointwise equality.

\subsection{Row 145: one fixed logarithm on the domain}
A complex reciprocal power needs a logarithm choice. The source makes it part of the data:
\excerpt{AnalyticReciprocalLaws}{fixed_log}
\excerpt{AnalyticReciprocalLaws}{log_power}
For one chosen $L$, the parameter law follows directly from the exponential addition identity.
\excerpt{AnalyticReciprocalLaws}{log_addition}
Differentiation requires an open domain and a holomorphic logarithm of the actual denominator.
\excerpt{AnalyticReciprocalLaws}{log_derivative}
The source also proves the differential Bell/Riccati identity at points where the denominator is nonzero. These statements close results 7.4, 8.4, and 8.5 alongside their formal coefficient laws. They do not choose unrelated logarithms for different parameters and then assert a branch-independent multiplicative law.

\subsection{Row 146: ordinary derivatives and Hasse derivatives}
The formal differential operator acts on polynomials, so its apparently infinite series truncates. There are two encodings: ordinary iterated derivatives and factorial-normalized Hasse derivatives. The bridge is explicit:
\excerpt{OperatorEncodingBridge}{operator_bridge}
The proof uses the identity relating the $q$th Hasse derivative to the ordinary $q$th derivative after multiplication by $q!$. The source then proves the forward action on the named residue Appell polynomial and its inverse.
\excerpt{OperatorEncodingBridge}{operator_inverse}
The inverse parameter is $-r$. The expansion of the residue series changes the differential powers by multiples of $N$. The polynomial degree makes the action finite, so the result does not require a Banach-space operator convergence theorem.

\subsection{Row 147: the order-one tangent boundary}
The order-one denominator is identified exactly:
\excerpt{TangentBoundary}{order_one}
It follows that the transfer tangent is one and its coefficients have the corresponding degenerate form. This boundary case is joined to the earlier $N\geq2$ convergence theorem to obtain every positive order.
\excerpt{TangentBoundary}{all_orders}
Notice the original-variable exponent $Nn+N-1$ and factorial $(Nn+N-1)!$. The nonvanishing hypothesis is stated in the reduced-variable ball, while the evaluation point is tested through $z^N$. These are explicit normalizations, not cosmetic choices.

\subsection{Row 148: literal compact reciprocal polynomials}
A recurrence determines low-order coefficients, but the manuscript displays compact binomial/factorial expressions. This module proves the displayed forms over every characteristic-zero field. For the second coefficient:
\excerpt{CompactReciprocal}{compact_two}
The third coefficient is likewise identified by its literal formula. The proof moves between ordinary coefficients and factorial normalization, justifies the nonzero factorial denominators, and uses the earlier recurrence. This closes the displayed-polynomial clause rather than relying on an informal simplification outside Lean.

\subsection{Row 149: the displayed Hessenberg determinant}
The determinant's matrix is defined entry by entry. Below the diagonal it contains the recurrence coefficients. Its superdiagonal contains negative integers.
\excerpt{NewtonDeterminant}{hessenberg}
The generic determinant theorem identifies the determinant with $n!b_n$. The reciprocal specialization is:
\excerpt{NewtonDeterminant}{determinant}
The finite-dimensional convention includes $n=0$, whose determinant and initial coefficient are one. The formula uses ordinary coefficients. Obtaining the manuscript's block coefficients requires their separate $(Nn)!$ conversion. The negative superdiagonal and the recurrence index are part of the checked matrix, not inferred from a picture.

\subsection{Row 150: hypergeometric normalization and summation}
The literal denominator parameters are generated by removing the unique residue parameter equal to one.
\excerpt{HypergeometricAnalytic}{hyper_parameters}
There are $N-1$ remaining parameters, all with positive real part. The source first relates factorials to residue Pochhammer products. The analytic closure then converts the pinned Mathlib regularized hypergeometric series to the ordinary normalization using Gamma factors and proves its sum.
\excerpt{HypergeometricAnalytic}{hyper_sum}
Finally the entire transfer section is identified at arbitrary complex arguments.
\excerpt{HypergeometricAnalytic}{hyper_bridge}
The proof reindexes the original Taylor sum by $Nq+j$, proves off-residue terms vanish, and matches the factorial/Pochhammer coefficients. The argument retains the scale $(\lambda z)^N/N^N$ and the prefactor $(\lambda z)^j/j!$. For $N=1$, the erased parameter multiset is empty and the same theorem applies.

\subsection{Row 151: convergent quotient-algebra exponentials}
The actual exponential series is defined in a topological algebra.
\excerpt{QuotientExponential}{quotient_exp}
Its definition alone does not assert summability. The next theorem supplies convergence for the cyclic relation under the surrounding module's topological linear-space assumptions.
\excerpt{QuotientExponential}{quotient_sum}
The proof groups powers by residue and uses the scalar \code{HasSum} theorem for each coordinate. The quotient algebras receive their finite-dimensional coefficient norm topologies and compatible normed-space instances. The signed specialization is:
\excerpt{QuotientExponential}{quotient_identity}
This is the final analytic link for result 6.2. The inherited formal power-series identity, the finite algebra relation, and the pointwise convergent exponential all retain the same generator and normalization.
